\section{Model danych}

\subsection{Schemat relacyjny}

PostgreSQL i MariaDB używają tego samego schematu z~12 tabel: \texttt{kategorie},
\texttt{organizatorzy}, \texttt{obiekty}, \texttt{sektory}, \texttt{miejsca},
\texttt{wydarzenia}, \texttt{cenniki}, \texttt{klienci}, \texttt{kody\_rabatowe},
\texttt{rezerwacje}, \texttt{bilety} i \texttt{platnosci}.

Bilet oznacza jedno miejsce oferowane na jedno wydarzenie. Jest pulą
sprzedażową: przy imporcie powstaje bilet dla każdego miejsca obiektu, a jego
status przechodzi z~\texttt{wolny} do \texttt{zablokowany},
\texttt{zarezerwowany} albo \texttt{sprzedany}. Dzięki temu rezerwacja miejsca
jest warunkową aktualizacją jednego wiersza
(\texttt{UPDATE bilety SET status = 'zarezerwowany' ... WHERE id = ? AND status = 'wolny'}).
Ta sama operacja istnieje w~każdym z~czterech modeli, co pozwala porównać
silniki na identycznym problemie.

\begin{figure}[H]
\centering
\begin{tikzpicture}[
  tbl/.style={draw, rounded corners=2pt, fill=gray!8, font=\footnotesize\ttfamily, minimum width=2.3cm, minimum height=0.6cm},
  rel/.style={-{Stealth[length=2mm]}, gray!70!black}]
  \node[tbl] (obj) {obiekty};
  \node[tbl, right=1.1cm of obj] (sek) {sektory};
  \node[tbl, right=1.1cm of sek] (mie) {miejsca};
  \node[tbl, below=0.9cm of obj] (wyd) {wydarzenia};
  \node[tbl, left=0.8cm of wyd] (kat) {kategorie};
  \node[tbl, below=0.9cm of kat] (org) {organizatorzy};
  \node[tbl, below=0.9cm of wyd] (cen) {cenniki};
  \node[tbl, below=0.9cm of mie] (bil) {bilety};
  \node[tbl, below=0.9cm of bil] (rez) {rezerwacje};
  \node[tbl, right=0.8cm of rez] (pla) {platnosci};
  \node[tbl, below=0.9cm of rez] (kli) {klienci};
  \node[tbl, left=0.8cm of kli] (kod) {kody\_rabatowe};
  \draw[rel] (sek) -- (obj); \draw[rel] (mie) -- (sek);
  \draw[rel] (wyd) -- (obj); \draw[rel] (wyd) -- (kat); \draw[rel] (wyd) -- (org);
  \draw[rel] (cen) -- (wyd); \draw[rel] (cen.east) -| (sek.south);
  \draw[rel] (bil) -- (mie); \draw[rel] (bil.west) -- ++(-1.2,0) |- (wyd.east);
  \draw[rel] (bil) -- (rez); \draw[rel] (pla) -- (rez);
  \draw[rel] (rez) -- (kli); \draw[rel] (rez.west) -| (kod.north);
\end{tikzpicture}
\caption{Schemat relacyjny (strzałka: kolumna odwołująca się do klucza głównego).}
\end{figure}

Klucze obce zostały celowo pominięte. InnoDB automatycznie tworzy indeks dla
każdego klucza obcego, więc w~wariancie ,,bez indeksów'' MariaDB miałaby
kilka indeksów, których PostgreSQL by nie miał. Spójność referencyjną
zapewnia aplikacja, a~sprawdzają ją testy (rozdz.~\ref{sec:testy}).

Wariant ,,z~indeksami'' dodaje w~obu bazach te same 15 indeksów drugorzędnych.
Wśród nich są \texttt{bilety(wydarzenie\_id, status)},
\texttt{bilety(rezerwacja\_id)}, \texttt{rezerwacje(klient\_id, utworzono)}
i \texttt{rezerwacje(status, wygasa\_o)}. Na danych JSON indeksujemy pole
\texttt{metadane->>'gatunek'}: w~PostgreSQL przez indeks na wyrażeniu
i~indeks GIN \texttt{jsonb\_path\_ops}, w~MariaDB przez kolumnę wirtualną
\texttt{gatunek} z~indeksem B-drzewa.

\subsection{Model dokumentowy (MongoDB)}

\begin{itemize}
  \item \texttt{wydarzenia}: wydarzenie z~osadzonym cennikiem (tablica
  \texttt{cennik}) oraz zdenormalizowanymi polami \texttt{miasto},
  \texttt{obiekt\_nazwa} i \texttt{kategoria\_nazwa}, żeby wyszukiwanie po
  mieście nie wymagało \texttt{\$lookup};
  \item \texttt{obiekty}: obiekt z~osadzonymi sektorami;
  \item \texttt{rezerwacje}: rezerwacja z~osadzonymi migawkami biletów
  (\texttt{bilety}: id, wydarzenie, sektor, rząd, numer, cena) i~osadzoną
  płatnością (\texttt{platnosc});
  \item \texttt{bilety}: pula sprzedażowa jako osobna kolekcja, ponieważ to ją
  modyfikują równolegli klienci; osadzenie 2--5 tys. biletów w~dokumencie
  wydarzenia zamieniłoby rezerwację różnych miejsc w~konflikt o~jeden dokument;
  \item \texttt{klienci}, \texttt{organizatorzy}, \texttt{kategorie},
  \texttt{kody\_rabatowe}.
\end{itemize}

MongoDB działa jako jednowęzłowy \emph{replica set} \texttt{rs0}, bo
transakcje wielodokumentowe wymagają replikacji. Operacje zmieniające kilka
dokumentów (np. rezerwacja z~biletami, usunięcie konta klienta) wykonujemy
w~transakcjach, tak jak w~bazach relacyjnych.

\subsection{Model klucz--wartość (Valkey)}

\begin{center}
\small
\begin{tabular}{llp{8.2cm}}
\toprule
klucz & typ & zawartość \\
\midrule
\texttt{w:\{id\}} & hash & wydarzenie: nazwa, miasto, data, organizator, metadane (JSON), cennik (JSON) \\
\texttt{wb:\{wydarzenie\}} & hash & id biletu $\to$ JSON biletu (miejsce, sektor, cena, status, rezerwacja) \\
\texttt{wf:\{wydarzenie\}} & set & identyfikatory wolnych biletów \\
\texttt{r:\{id\}} & string & rezerwacja jako JSON z~biletami i~płatnością \\
\texttt{k:\{id\}} & hash & klient, preferencje jako JSON \\
\texttt{blok:\{bilet\}} & string z~TTL & czasowa blokada miejsca \\
\texttt{sprzedaz}, \texttt{przychod} & zset, hash & liczniki sprzedaży i~przychodu per wydarzenie \\
\texttt{ix:*} & set, zset & struktury indeksujące, tylko w~wariancie z~indeksami \\
\bottomrule
\end{tabular}
\end{center}

Valkey nie ma planera zapytań ani indeksów drugorzędnych. W~wariancie
,,z~indeksami'' aplikacja utrzymuje dodatkowe zbiory posortowane, np.
\texttt{ix:k:\{klient\}} (rezerwacje klienta według daty),
\texttt{ix:miasto:\{miasto\}} (wydarzenia według daty) czy \texttt{ix:oczek}
(rezerwacje według terminu wygaśnięcia). Bez nich zapytanie po atrybucie
wymaga \texttt{SCAN} po całej przestrzeni kluczy i~dekodowania każdego
dokumentu. Operacje zmieniające kilka kluczy są atomowe dzięki skryptom Lua.
Liczniki sprzedaży są częścią modelu bazowego, bo w~tej domenie są naturalną
strukturą, a~nie indeksem.

\subsection{Dane i rozmiary zbiorów}

Dane tworzy deterministyczny generator. Każda encja wynika z~trójki (ziarno,
rodzaj, identyfikator), więc każdy serwis generuje dane u~siebie i~wszystkie
cztery bazy dostają identyczne wiersze. Po imporcie serwis zwraca sumę
kontrolną FNV-64 wygenerowanych danych. Raport porównuje ją między bazami.

\begin{table}[H]
\centering
\caption{Rozmiary zbiorów (ziarno 42).}
\begin{tabular}{rrrrrr}
\toprule
bilety & wydarzenia & obiekty & miejsca & klienci & rezerwacje \\
\midrule
\num{500000} & 244 & 8 & \num{16684} & \num{62500} & \num{82807} \\
\num{1000000} & 476 & 16 & \num{33922} & \num{125000} & \num{170396} \\
\num{10000000} & \num{4156} & 166 & \num{396795} & \num{1250000} & \num{1674389} \\
\bottomrule
\end{tabular}
\end{table}

Około 60\% biletów jest sprzedanych albo zarezerwowanych, rezerwacja obejmuje
od 1 do 6 sąsiednich miejsc, 85\% rezerwacji jest opłaconych, a~10\% ma kod
rabatowy. Metadane wydarzeń (gatunek, wykonawcy, limit wiekowy, tagi),
biletów, klientów, rezerwacji i~płatności są zapisane jako JSON. Cztery
wydarzenia w~każdym zbiorze nie mają sprzedaży (,,piaskownica''). Korzystają
z~nich scenariusze rezerwujące miejsca i~test współbieżności, więc nigdy nie
kolidują z~danymi ani ze sobą nawzajem.
